\section{Introduction}

The blood--brain barrier is a selective endothelial boundary that excludes most circulating compounds from the central nervous system. A drug intended to act in the brain is useless if it cannot cross it, and one intended to act elsewhere is a liability if it does. Measuring penetration requires animal or tissue assays, so it is screened computationally first, when a library holds millions of structures and the laboratory can handle hundreds.

Two representations dominate that screen. A \emph{fingerprint} hashes the substructures present in a molecule into a fixed bit vector that a linear model can classify, and cannot adapt to the task; a \emph{graph neural network} treats the molecule as a graph of atoms and bonds and learns its own atom representations by message passing. The second is the more expressive, and the question here is whether that expressiveness pays on a dataset of realistic size.

The dataset is BBBP from MoleculeNet \parencite{MoleculeNet}, assembled by Martins et al.\ \parencite{Martins}: 2039 molecules loaded as graphs with a binary penetration label, 76.5\,\% of them positive. Two choices decide what the numbers mean. A random split lets a test molecule share its ring system with a training molecule, rewarding recognition of a chemotype already measured, whereas a scaffold split forbids it and so asks the question deployment asks; and a 76\,\% positive base rate makes accuracy weak, since predicting ``permeable'' for everything already scores 0.763. This report audits the dataset, builds both splits, trains a fingerprint baseline and an edge-aware graph network under each, evaluates both on sealed test sets, and analyses the network's most confident errors against the training distribution.
